
%---------------------------------------------------------------
\chapter{Introduction}

% uncomment the following line to create an unnumbered chapter
%\chapter*{Introduction}\addcontentsline{toc}{chapter}{Introduction}\markboth{Introduction}{Introduction}
%---------------------------------------------------------------
\setcounter{page}{1}

% The following environment can be used as a mini-introduction for a chapter. Use that any way it pleases you (or comment it out). It can contain, for instance, a summary of the chapter. Or, there can be a quotation.



\begin{markdown}

This thesis is primarily focused on the practical application of modern natural language processing methods for fake news classification. It builds upon the bachelor thesis by Jan Flajžík, Classification of Fake News in The Media/Social Media Ecosystem \cite{flajzik_classification_2024}, which explored the use of classical machine learning models for predicting whether a news article is fake or reliable. The present work extends this research by employing more advanced transformer-based architectures.

The theoretical part of the thesis is limited mainly to the description of the selected models and their underlying principles. Rather than providing a deep theoretical analysis of the architectures, this work approaches the problem from a practical perspective by applying advanced models, evaluating their performance, and comparing the obtained results with those reported in previous work.

The datasets used in this work remain identical to those utilized in the previous work \cite{flajzik_classification_2024}. The primary difference lies in the use of computationally more demanding transformer models and in the integration of modern tools for experiment tracking and analysis, specifically the MLFlow framework.

The motivation for this research is well described in the previous work \cite{flajzik_classification_2024} and remains highly relevant today. The rapid spread of manipulative and misleading information through online platforms represents a significant societal issue. Due to the enormous volume of digital content published daily, manual verification of information is economically and practically infeasible. Consequently, automated fake news detection using machine learning techniques has become an important research area.

The following chapters introduce the selected transformer models, describe the implementation and experimental methodology, present the achieved results, and compare them with those obtained in previous research. Finally, the thesis discusses the limitations of the proposed approach and outlines possible directions for future work.
      
\end{markdown}

%---------------------------------------------------------------
\section{Goals}
%---------------------------------------------------------------

\begin{markdown}

The primary objective of this thesis is to extend previous research in fake news classification by integrating advanced natural language processing architectures, particularly transformer-based models such as BERT, GPT, and related approaches. The work focuses on evaluating whether these modern architectures can improve classification performance compared to the classical machine learning methods used in earlier studies.

The thesis also examines different text preprocessing strategies and their influence on model performance. Experimental evaluation is conducted on the same datasets as those used in the previous work \cite{flajzik_classification_2024}, enabling a direct comparison between traditional and transformer-based approaches.
    
The main goals of this thesis are as follows:

1. Investigate state-of-the-art methods for fake news classification.
2. Implement and evaluate at least three advanced transformer-based models, such as BERT, GPT, or related architectures.
3. Conduct experiments using different preprocessing techniques and compare the obtained results with those reported in previous work.
4. Analyze the strengths and limitations of the selected models and discuss possible improvements for future research.

\end{markdown}